% GENERATED FILE. Edit canonical lessons, then run npm run generate:books.
% canonical-source-hash: fnv1a64:3f871cf417b24d5a
% canonical-lessons: MR-C536-gayan, MR-C536-vadak, MR-C536-sur-note, MR-C536-gungunne, MR-C536-riyaj-karne

\chapter{Music Making}
\label{ch:mr-music-making}
\hlchaptermodality{536}

\begin{chapteropening}
\textbf{By the end of this chapter:} \emph{I can talk about singing, notes and practising music.}

Complete the last lesson of chapter 536: I can talk about singing, notes and practising music.
\end{chapteropening}

\section[\texorpdfstring{\mr{गायन} (gāyan)}{गायन (gāyan)}]{\mr{गायन} (gāyan) — singing as an art}

\label{lesson:MR-C536-gayan}



\begin{quote}
Before the new one: say the Marathi for a paintbrush, then the Marathi for a shade.
\end{quote}

\subsection*{You'll want to know: \mr{गायन}}
\textbf{\mr{गायन}} — \emph{gāyan} — \textquotedblleft{}singing as an art\textquotedblright{}.

\begin{quote}
\mr{संध्याकाळी} \mr{शाळेत} \mr{गायनाचा} \mr{कार्यक्रम} \mr{आहे}.
\end{quote}

\begin{quote}
\emph{sandhyākāḷī śāḷet gāyanācā kāryakram āhe.}
\end{quote}

\begin{quote}
There is a singing programme at school this evening.
\end{quote}

Neuter.

\begin{cousinweb}
\textbf{\mr{गायन}} (\emph{gāyan}) is a tatsama, taken from Sanskrit \emph{gāyana}, \textquotedblleft{}a singer, singing\textquotedblright{} (Monier-Williams).
\end{cousinweb}

\begin{grammarlens}[title={What you've built}]
One more word for describing what you lived and felt.
\end{grammarlens}

\subsection*{Your turn}
\begin{itemize}
\raggedright
  \item \emph{Say it:} \emph{gāyan}
  \item \emph{Say it:} \emph{gāyan}, once more
  \item \emph{Say it:} \emph{sandhyākāḷī śāḷet gāyanācā kāryakram āhe.} — There is a singing programme at school this evening.
  \item \emph{Recall:} say the Marathi for a paintbrush, then the Marathi for a shade, then say \emph{gāyan} again
\end{itemize}

\begin{tcolorbox}[breakable,colback=teal!4,colframe=teal!35!black,title={Before you move on}]
What does \mr{गायन} mean? (\textquotedblleft{}Singing as an art\textquotedblright{}.) Say it once more.
\end{tcolorbox}
\section[\texorpdfstring{\mr{वादक} (vādak)}{वादक (vādak)}]{\mr{वादक} (vādak) — an instrumentalist}

\label{lesson:MR-C536-vadak}



\begin{quote}
Before the new one: say the Marathi for a shade, then the Marathi for singing as an art.
\end{quote}

\subsection*{You'll want to know: \mr{वादक}}
\textbf{\mr{वादक}} — \emph{vādak} — \textquotedblleft{}an instrumentalist\textquotedblright{}.

\begin{quote}
\mr{तबला} \mr{वादकाने} \mr{गायकाला} \mr{सुंदर} \mr{साथ} \mr{दिली}.
\end{quote}

\begin{quote}
\emph{tablā vādakāne gāyakālā sundar sāth dilī.}
\end{quote}

\begin{quote}
The tabla player accompanied the singer beautifully.
\end{quote}

Masculine.

\mr{साथ} \mr{देणे}, to accompany.

\begin{cousinweb}
\textbf{\mr{वादक}} (\emph{vādak}) is a tatsama, taken from Sanskrit \emph{vādaka}, \textquotedblleft{}one who makes sound, a player\textquotedblright{} (Monier-Williams).
\end{cousinweb}

\begin{grammarlens}[title={What you've built}]
One more word for describing what you lived and felt.
\end{grammarlens}

\subsection*{Your turn}
\begin{itemize}
\raggedright
  \item \emph{Say it:} \emph{vādak}
  \item \emph{Say it:} \emph{vādak}, once more
  \item \emph{Say it:} \emph{tablā vādakāne gāyakālā sundar sāth dilī.} — The tabla player accompanied the singer beautifully.
  \item \emph{Recall:} say the Marathi for a shade, then the Marathi for singing as an art, then say \emph{vādak} again
\end{itemize}

\begin{tcolorbox}[breakable,colback=teal!4,colframe=teal!35!black,title={Before you move on}]
What does \mr{वादक} mean? (\textquotedblleft{}An instrumentalist\textquotedblright{}.) Say it once more.
\end{tcolorbox}
\section[\texorpdfstring{\mr{सूर} (sūr)}{सूर (sūr)}]{\mr{सूर} (sūr) — a musical note; a tune}

\label{lesson:MR-C536-sur-note}



\begin{quote}
Before the new one: say the Marathi for singing as an art, then the Marathi for an instrumentalist.
\end{quote}

\subsection*{You'll want to know: \mr{सूर}}
\textbf{\mr{सूर}} — \emph{sūr} — \textquotedblleft{}a musical note; a tune\textquotedblright{}.

\begin{quote}
\mr{गायिकेचा} \mr{पहिलाच} \mr{सूर} \mr{खूप} \mr{गोड} \mr{होता}.
\end{quote}

\begin{quote}
\emph{gāyikecā pahilāc sūr khūp goḍ hotā.}
\end{quote}

\begin{quote}
The singer's very first note was very sweet.
\end{quote}

Masculine.

\begin{cousinweb}
\textbf{\mr{सूर}} (\emph{sūr}) is a tadbhava: CDIAL lists Marathi \emph{sūr}, \textquotedblleft{}musical note\textquotedblright{}, under Sanskrit \emph{sura}, \textquotedblleft{}sound\textquotedblright{}, a variant of \emph{svara} (CDIAL 13498).
\end{cousinweb}

\begin{grammarlens}[title={What you've built}]
One more word for describing what you lived and felt.
\end{grammarlens}

\subsection*{Your turn}
\begin{itemize}
\raggedright
  \item \emph{Say it:} \emph{sūr}
  \item \emph{Say it:} \emph{sūr}, once more
  \item \emph{Say it:} \emph{gāyikecā pahilāc sūr khūp goḍ hotā.} — The singer's very first note was very sweet.
  \item \emph{Recall:} say the Marathi for singing as an art, then the Marathi for an instrumentalist, then say \emph{sūr} again
\end{itemize}

\begin{tcolorbox}[breakable,colback=teal!4,colframe=teal!35!black,title={Before you move on}]
What does \mr{सूर} mean? (\textquotedblleft{}A musical note; a tune\textquotedblright{}.) Say it once more.
\end{tcolorbox}
\section[\texorpdfstring{\mr{गुणगुणणे} (guṇguṇṇe)}{गुणगुणणे (guṇguṇṇe)}]{\mr{गुणगुणणे} (guṇguṇṇe) — to hum}

\label{lesson:MR-C536-gungunne}



\begin{quote}
Before the new one: say the Marathi for an instrumentalist, then the Marathi for a musical note.
\end{quote}

\subsection*{You'll want to know: \mr{गुणगुणणे}}
\textbf{\mr{गुणगुणणे}} — \emph{guṇguṇṇe} — \textquotedblleft{}to hum\textquotedblright{}.

\begin{quote}
\mr{स्वयंपाक} \mr{करताना} \mr{आई} \mr{एक} \mr{जुनं} \mr{गाणं} \mr{गुणगुणत} \mr{होती}.
\end{quote}

\begin{quote}
\emph{svayampāk kartānā āī ek junaṁ gāṇaṁ guṇguṇat hotī.}
\end{quote}

\begin{quote}
While cooking, Mother was humming an old song.
\end{quote}

Past, with -\mr{ने} (he …): \mr{त्याने} \mr{गाणं} \mr{गुणगुणलं}, \emph{tyāne gāṇaṁ guṇguṇlaṁ}.

\begin{cousinweb}
\textbf{\mr{गुणगुणणे}} (\emph{guṇguṇṇe}) is an inherited echo word: CDIAL sets Hindi \emph{gungunānā} and the Marathi murmuring words under \emph{*ganagana}, \emph{*gunaguna}, \textquotedblleft{}murmur\textquotedblright{} (CDIAL 4013).
\end{cousinweb}

\begin{grammarlens}[title={What you've built}]
One more word for describing what you lived and felt.
\end{grammarlens}

\subsection*{Your turn}
\begin{itemize}
\raggedright
  \item \emph{Say it:} \emph{guṇguṇṇe}
  \item \emph{Say it:} \emph{guṇguṇṇe}, once more
  \item \emph{Say it:} \emph{svayampāk kartānā āī ek junaṁ gāṇaṁ guṇguṇat hotī.} — While cooking, Mother was humming an old song.
  \item \emph{Recall:} say the Marathi for an instrumentalist, then the Marathi for a musical note, then say \emph{guṇguṇṇe} again
\end{itemize}

\begin{tcolorbox}[breakable,colback=teal!4,colframe=teal!35!black,title={Before you move on}]
What does \mr{गुणगुणणे} mean? (\textquotedblleft{}To hum\textquotedblright{}.) Say it once more.
\end{tcolorbox}
\section[\texorpdfstring{\mr{रियाज} \mr{करणे} (riyāj karṇe)}{रियाज करणे (riyāj karṇe)}]{\mr{रियाज} \mr{करणे} (riyāj karṇe) — to practise music}

\label{lesson:MR-C536-riyaj-karne}



\begin{quote}
Before the new one: say the Marathi for a musical note, then the Marathi for to hum.
\end{quote}

\subsection*{You'll want to know: \mr{रियाज} \mr{करणे}}
\textbf{\mr{रियाज} \mr{करणे}} — \emph{riyāj karṇe} — \textquotedblleft{}to practise music\textquotedblright{}.

\begin{quote}
\mr{गायक} \mr{रोज} \mr{पहाटे} \mr{रियाज} \mr{करतो}.
\end{quote}

\begin{quote}
\emph{gāyak roj pahāṭe riyāj karto.}
\end{quote}

\begin{quote}
The singer practises every day at dawn.
\end{quote}

Past, with -\mr{ने} (he …): \mr{त्याने} \mr{रियाज} \mr{केला}, \emph{tyāne riyāj kelā}.

\begin{cousinweb}
\textbf{\mr{रियाज} \mr{करणे}} (\emph{riyāj karṇe}) is \mr{रियाज}, a loan through Persian usually traced to Arabic \emph{riyāḍa}, \textquotedblleft{}training\textquotedblright{}, with \mr{करणे}.
\end{cousinweb}

\begin{grammarlens}[title={What you've built}]
One more word for describing what you lived and felt.
\end{grammarlens}

\subsection*{Your turn}
\begin{itemize}
\raggedright
  \item \emph{Say it:} \emph{riyāj karṇe}
  \item \emph{Say it:} \emph{riyāj karṇe}, once more
  \item \emph{Say it:} \emph{gāyak roj pahāṭe riyāj karto.} — The singer practises every day at dawn.
  \item \emph{Recall:} say the Marathi for a musical note, then the Marathi for to hum, then say \emph{riyāj karṇe} again
\end{itemize}

\begin{tcolorbox}[breakable,colback=teal!4,colframe=teal!35!black,title={Before you move on}]
What does \mr{रियाज} \mr{करणे} mean? (\textquotedblleft{}To practise music\textquotedblright{}.) Say it once more.
\end{tcolorbox}
